% 10장 논의. framing.md 2절 · 0022 · 0031

\section{논의}

\subsection{명분과 작동의 어긋남}

이 연구의 뼈대를 한 문장으로 쓰면 이렇다.

\begin{quote}
\textbf{토지거래허가제의 명분은 가격 안정인데 실제로 작동한 것은 거래
동결이었다.} 거래는 확실하게 끊었지만 가격을 잡았다는 증거는 약하다.
\end{quote}

이 어긋남은 두 숫자의 대비에 있다. 지정 직후 거래량은 반사실 대비
$-71 \sim -80\%$이고, 설계 조합 192가지 전부에서 음수이며, 사전 추세를 실제
관측값의 1.6배까지 열어 두어도 살아남고, 서로 다른 네 자치구에서 열 달 뒤에
일어난 별개 사건에서도 같은 크기다. 반면 가격은 평행추세를 가정해야 $-4 \sim
-7\%$이고 그 가정을 조금만 풀면 0과 구분되지 않으며 합성대조군으로는 확인되지
않는다.

\textbf{이것은 제도가 실패했다는 말이 아니다.} 제도가 직접 건드리는 것은 거래가
성립하는 절차이지 가격이 아니므로, 거래가 멈춘 것은 제도의 설계대로 작동한
결과다. 말할 수 있는 것은 \textbf{그 작동이 명분까지 닿았다는 증거가 이 자료에서
나오지 않는다}는 것이다.

\subsection{설계 세부가 효과를 가른다 --- 정책 시사점}

\textbf{이 연구의 가장 실용적인 결론은 6장에 있다.}

2020년 지정의 허가 대상은 주거지역 대지 \textbf{18\,m$^2$ 초과}였다. 그 선 아래의
소형 아파트는 구역 안에 있어도 허가 없이 거래할 수 있었고, 실제로 거래가 줄지
않았다 --- 줄지 않은 정도가 아니라 \textbf{통제군 대비 두 배 이상으로 늘었다}
($+125.7\%$). 규제가 수요를 없앤 것이 아니라 \textbf{면제 구간으로 밀어냈다.}

2022년 6월 23일 그 기준이 \textbf{6\,m$^2$}로 좁아졌다. 그러자 \emph{그때} 소형
거래가 약 3분의 1 줄었고($-33.2\%$), 기준을 바꾸지 않은 2021년 재지정에서는
움직이지 않았다($-8.1\%$).

\begin{quote}
\textbf{규제를 ``거느냐 마느냐''가 아니라 ``어떻게 거느냐''가 효과를 정한다.}
가격 차등(대출 규제)보다 설계 세부(허가 대상의 면적 기준)가 규제의 실효를 갈랐다.
\end{quote}

이 결론에는 두 가지 함의가 있다.

\textbf{첫째, 면적 기준이 사실상의 회피 통로였다.} 규제를 피할 수 있는 물건은
거래가 유지됐고 통로가 닫히자 줄었다. 제도 설계자가 2년 뒤 기준을 6\,m$^2$로
고친 것은 규제 당국도 결국 같은 것을 알아챘음을 시사한다.

\textbf{둘째, 정책 평가가 평균만 보면 이 사실이 보이지 않는다.} 동 평균 안에서
대형의 $-95\%$와 소형의 $+126\%$가 상쇄되어 하나의 숫자로 뭉개진다. 거래마다
전용면적이 붙은 미시 자료가 아니면 \textbf{애초에 볼 수 없다.} 선행연구 세 편이
동 단위 집계로 ``60\% 줄었다''에서 멈춘 이유가 여기 있다.

\subsection{거래 동결 자체가 비용이다}

``거래가 줄었다''를 중립적인 사실로 읽으면 안 된다. 거래가 80\% 끊기면 이사가
필요한 사람도 옮기지 못하고, 시장이 값을 찾는 신호도 약해진다. \textbf{효과의
크기와 지속을 재는 일은 그 비용을 재는 일이기도 하다.}

그래서 ``언제까지 갔나''가 중요한 질문이었고, \textbf{이 연구는 그 질문에
답하지 못했다}(다음 절).

\subsection{한계}

숨기지 않고 적는다.

\textbf{일반화의 범위.} 처치가 두 사건·아홉 개 동이고 \textbf{두 집단 모두 서울
고가 재건축 지역}이다. 2021년 B그룹에서 같은 크기의 효과가 나왔고 두 사건의
차이가 0과 구분되지 않았지만(9.6절), 완전히 다른 유형의 지역은 아니다.
\textbf{``서울 고가 재건축 지역에서는 일관된다''까지가 정확하고}, 지방이나 다른
유형의 지역에는 말할 수 없다.

\textbf{지속 기간을 말할 수 없다.} 사후 $k \geq 4$는 기준 분기 선택에 따라
$-20 \sim +16\%$로 흔들리고, 24개월 평균의 붕괴값이 0.05이며, 합성대조군에서
사후 27개월 평균은 152개 중 22위($p=0.145$)다. 세 진단이 같은 방향을 가리킨다.
\textbf{일시 충격인지 지속적 동결인지는 이 설계로 답할 수 없다.}

\textbf{가격의 크기를 말할 수 없다.} 방향까지만 쓴다(8.6절).

\textbf{부과와 해제의 대칭성을 가릴 수 없다.} $\text{DID}_M$에서 부과 쪽은 구간이
0을 품지 않고 해제 쪽은 품는다(9.4절).

\textbf{대출의 몫을 가릴 수 없다.} 면적 기준이 주된 원인이라는 것이 대출의 몫이
0이라는 뜻은 아니다. 소형을 뺀 뒤에도 15억 경계의 차이가 점추정으로 남는다
(무작위화 $p=0.098$). 폐지를 사건으로 걸어 한 번 더 시도했으나 금리 정점 통과와
1·3 대책에서 가를 수 없었다(6.9절).

\textbf{대지지분이 자료에 없다.} 허가 기준은 대지지분인데 실거래 자료에 없어
전용면적으로 대신했다. 2021년 B그룹에서 더 부정확하다 --- 1970$\sim$80년대 재건축
대상 단지라 용적률이 낮아 같은 전용면적이라도 대지지분이 클 수 있다. B그룹 소형
검정이 방향만 맞고 세워지지 않은 이유의 하나다.

\textbf{거래 주체를 보지 못했다.} 법인 매수와 직거래 비중은 규제 회피 경로를
직접 가리키는 변수인데, \texttt{dealingGbn}이 2021년 말, \texttt{slerGbn}·
\texttt{buyerGbn}이 2024년부터 채워져 \textbf{2020·2021 사건에는 원천적으로 쓸 수
없다}(2020년은 12개월 전부 결측). 2025년 사건에만 적용 가능하다.

\textbf{취소 플래그가 2020년 2월부터다.} 사건연구 사전 구간의 대부분에서 취소를
뺄 수 없고 사후에서는 뺀다. 재 보니 그 비대칭이 주 결과를 0.3\,\%p 움직인다
(9.7절).

\textbf{저장소가 무겁다.} 그림을 매번 커밋해 \texttt{.git}이 936\,MB다. 이력을
다시 쓰는 일이라 이번 학기에는 고치지 않았다.
